\section{Introduction}
% =====================================================================

Galaxy surveys such as DESI map the large-scale structure of the Universe over volumes large enough that statistical errors are no longer the main limitation. How the galaxies are distributed tells us how structures grew, how much matter there is, and what dark energy and neutrinos do. The question is how much of this information our statistics actually extract.

Most analyses compress the galaxy field into its two-point statistics, the correlation function or the power spectrum. For a Gaussian field this loses nothing, and on large scales it works very well. But gravity makes the late-time Universe strongly non-Gaussian: matter collapses into halos, gathers in filaments and nodes, and empties voids, and peculiar velocities distort the picture in redshift space. The two-point function does not see all of this. Statistics built on the cosmic web recover part of the missing information: \cite{Sunseri} showed with Quijote that splitting the density field into nodes, filaments, walls and voids and combining the power spectra of these environments tightens the constraints substantially, most of all on the neutrino mass.

For galaxies, two families of environment statistics have been developed for this purpose. Density-split clustering \cite{Paillas2023, Paillas2024} ranks random points by the smoothed galaxy density around them, splits them into quantiles, and measures the auto-correlation of each quantile and its cross-correlation with the galaxies; combined with a simulation-based emulator \cite{CuestaLazaro2024}, it improves $\Lambda$CDM constraints by factors of several over the two-point function. Void statistics, combined with galaxy clustering in a forward model of the galaxy--halo connection, also break degeneracies between cosmology and galaxy formation \cite{Salcedo}.

ASTRA \cite{ASTRA} takes a different route. It triangulates the galaxies together with a uniform random catalogue and gives every point an environment from the fraction of galaxies among its Delaunay neighbours: there is no smoothing scale, and the environment adapts to the local density. ASTRA can classify points into voids, sheets, filaments and knots, or rank them continuously.

\paragraph{Goal.} Our aim is to quantify how much ASTRA environment statistics improve the constraints on the cosmological parameters ($\Omega_m$ through $\omega_{\rm cdm}$, $\sigma_8$, the dark-energy parameters $w_0$ and $w_a$, and ideally the neutrino mass) over the standard galaxy two-point function, and how they compare with density split when both are measured on the same mocks, with the same random points and the same estimators. To this end we build dedicated covariance and derivative datasets from suites of mock catalogues, which give Fisher forecasts of the parameter uncertainties, and we train a neural-network emulator of the statistics, which we use for an end-to-end inference test. The neutrino mass turned out to be out of reach with the available simulations (Section~\ref{sec:quijote}); all the other parameters are covered.

\paragraph{Data.} We use three simulation suites, each for what the others cannot provide (Section~\ref{sec:sims}): the AbacusSummit emulator suite of DESI-like galaxy mocks over 85 cosmologies, on which the measurements and the emulator are built; the official DESI mocks, whose 1786 independent boxes give a reliable covariance; and the Quijote simulations, whose paired simulations give clean derivatives. Getting reliable derivatives and covariances was the main difficulty of the work, and we describe how our choices were made and how the results evolved.

\paragraph{Outline.} Section~\ref{sec:env} presents the two environment estimators, Section~\ref{sec:clustering} the clustering measurements, Section~\ref{sec:framework} the covariance, the derivatives and the Fisher formalism, Section~\ref{sec:choices} the successive analysis choices, Section~\ref{sec:results} the results (DESI-like mocks, Quijote, ASTRA classes, inference), Section~\ref{sec:robustness} the robustness tests, and Section~\ref{sec:discussion} the discussion. The appendices explain how to reproduce the paper (\ref{app:repro}), give the mathematical details (\ref{app:maths}), and retrace the methodological path (\ref{app:path}).
